% answers_short.tex —— 1988 年·全国卷·文科 · 简短答案速查
% 由 tools/make_short_answers.py 自 answers.tex 的 \gksolhead/\gkanswer 提取，
% 并与 1988/1988_数学真题_全国文_答案要点.md 逐题比对后生成，勿手改（改源 MD 后重跑）。

\begin{gknotice}
\textbf{使用说明}：本节只给最终答案便于快速核对，完整推导见解析版。答案为编写者独立推导，\textbf{非官方标准答案}。
\end{gknotice}

\gkpart{选择题}
\begin{center}\begin{tabular}{cccccccccccccccc}
\toprule
题号 & 1 & 2 & 3 & 4 & 5 & 6 & 7 & 8 & 9 & 10 & 11 & 12 & 13 & 14 & 15 \\
\midrule
答案 & B & C & D & B & A & B & D & D & A & C & A & C & B & A & C \\
\bottomrule
\end{tabular}\end{center}

\gkpart{解答题}
\begin{itemize}[leftmargin=2.6em]
  \item[\textbf{16.}] 模为 $2$，辐角主值为 $-\dfrac\pi6$
  \item[\textbf{17.}] $x=-2$
  \item[\textbf{18.}] $\tan\dfrac\theta2=-3$
  \item[\textbf{19.}] $\dfrac{48\pi}{5}\text{ cm}^3$
  \item[\textbf{20.}] $3$
  \item[\textbf{21.}] 恒等式成立。
  \item[\textbf{22.}] $\sin\alpha=\dfrac1{\sqrt5}$
  \item[\textbf{23.}] $P=\left(\dfrac54,-\dfrac34\right)$
  \item[\textbf{24.}] $\displaystyle\left(-1,\frac{1-\sqrt5}{2}\right)\cup\left(1,\frac{1+\sqrt5}{2}\right)$
  \item[\textbf{25.}] $5m^2+m+2^{m+1}-2$
\end{itemize}
